\addcontentsline{toc}{subsubsection}{Solution 51. The Logarithm Hidden in a Chain}

\begin{dlsummary}[title={\bfseries Solution 51. The Logarithm Hidden in a Chain},breakable]

\begin{proof}
The unique solution is
\[
\boxed{f(x)=\log x\qquad(x>0),}
\]
where \(\log\) denotes the natural logarithm.

\paragraph{1. Find the normalisation and a pointwise bound.}

Setting \(x=y=1\), we get
\[
f(1)\le2f(1)\le0.
\]
Thus \(f(1)=0\).

Now set \(y=1\) in the upper inequality:
\[
\boxed{f(x)\le x-1.} \tag{1}
\]

\paragraph{2. Improve the upper bound by taking roots.}

Repeated application of the first inequality gives
\[
f(t^n)\le n f(t)
\]
for every positive integer \(n\).

Taking \(t=x^{1/n}\) and using (1),
\[
f(x)\le n f(x^{1/n})
\le n\bigl(x^{1/n}-1\bigr).
\]
The standard elementary limit
\[
\lim_{n\to\infty}n\bigl(x^{1/n}-1\bigr)=\log x
\]
therefore yields
\[
\boxed{f(x)\le\log x.} \tag{2}
\]

\paragraph{3. Use reciprocals to obtain the reverse bound.}

The first inequality at \((x,1/x)\) gives
\[
0=f(1)\le f(x)+f(1/x).
\]
Applying (2) at \(1/x\),
\[
f(x)\ge-f(1/x)\ge-\log(1/x)=\log x.
\]
Together with (2), this proves \(f(x)=\log x\).

\paragraph{4. Verify.}

The logarithm satisfies
\[
\log(xy)=\log x+\log y,
\]
and the standard inequality \(\log t\le t-1\) gives
\[
\log x+\log y\le x+y-2.
\]
Thus it satisfies the entire chain.
\end{proof}

\paragraph{Takeaways.}
\begin{itemize}
\item **Multiplicative subdivision sharpens bounds:** Replacing \(x\) by \(n\) identical factors improves \(f(x)\le x-1\) to the exact logarithmic bound.
\item **Reciprocals reverse one-sided information:** An upper bound at \(1/x\) becomes a lower bound at \(x\).
\item **Normalisation can be forced:** Neither \(f(1)=0\) nor a value at any other input needs to be supplied.
\item **No continuity of \(f\) is needed:** The limit concerns an explicit numerical expression, not function values approaching an input.
\end{itemize}

\end{dlsummary}
